\documentclass[a4paper,12pt]{article}

\usepackage[numbers]{natbib}
\usepackage{url}

\title{Explainable Machine Learning for Gallstone Disease}
\author{Matteo Soldini, Syed Sakibur Rahman, Irtaza Bin Ashar}

\begin{document}

\maketitle

\section{Introduction}
\label{sec:introduction}

We use the Gallstone dataset~\cite{gallstone_1150}.

The objective of this project is therefore twofold:
\begin{enumerate}
\item Develop a machine-learning model that can distinguish patients with and without gallstone disease.
\item Investigate which clinical characteristics contribute most to the model's predictions.
\end{enumerate}

\section{Methods}
\label{sec:methods}

\subsection{Dataset}
\label{ssec:dataset}

\subsubsection{Overview}
\label{ssec:overview}

The Gallstone dataset~\cite{gallstone_1150} was collected at the Internal
Medicine Outpatient Clinic of Ankara VM Medical Park Hospital between June 2022
and June 2023. It contains 319 records, one per patient, and 39 columns: the binary target
variable \emph{Gallstone Status} and 38 features. In the original coding, the
target takes the value 0 for patients with gallstones and 1 for patients
without.

The features fall into three groups:
\begin{itemize}
\item \textbf{Demographic:} age, gender, height, weight and body mass index (BMI);
\item \textbf{Bioimpedance:} body composition measurements, such as total,
  extracellular and intracellular body water, muscle and fat mass, protein
  content, visceral fat area and hepatic fat accumulation;
\item \textbf{Laboratory:} blood test results, including glucose, total
  cholesterol, LDL, HDL, triglycerides, liver enzymes (AST, ALT, ALP),
  creatinine, glomerular filtration rate, C-reactive protein (CRP), hemoglobin
  and vitamin D.
\end{itemize}
The dataset also records comorbidities (coronary artery disease,
hypothyroidism, hyperlipidemia and diabetes mellitus) as binary indicators,
together with the total number of comorbid conditions. Hepatic fat
accumulation is an ordinal grade from 0 (none) to 4 (very severe). All other
features are numerical.

The classes are nearly balanced: 161 patients (50.5\%) were diagnosed with
gallstone disease and 158 (49.5\%) were not
(Table~\ref{tab:class_balance}). The dataset contains no missing values and no
duplicate records.

\begin{table}[h]
\centering
\begin{tabular}{lrr}
\hline
Gallstone Status & Patients & Share \\
\hline
Present (0) & 161 & 50.5\% \\
Absent (1)  & 158 & 49.5\% \\
\hline
Total       & 319 & 100\% \\
\hline
\end{tabular}
\caption{Class distribution of the target variable.}
\label{tab:class_balance}
\end{table}

\subsubsection{Cleaning}
\label{ssec:cleaning}

\paragraph{Implausible Values}

To identify candidate outliers, we computed the mean and standard deviation of
each feature and flagged all values lying more than ten standard deviations
from the feature mean ($|z| > 10$). This conservative threshold targets only
extreme values that are likely to be data-entry errors rather than natural
variation. Five values in four records exceeded it
(Table~\ref{tab:outliers}). Each was then checked for clinical plausibility
and for consistency with related features in the same record.

\begin{table}[h]
\centering
\begin{tabular}{llrrrr}
\hline
Record & Feature & Value & Mean & Std & $z$ \\
\hline
40  & Glucose & 575.0 & 108.69 & 44.85 & 10.40 \\
239 & Obesity (\%) & 1954.0 & 35.85 & 109.80 & 17.47 \\
278 & Aspartat Aminotransferaz (AST) & 195.0 & 21.68 & 16.70 & 10.38 \\
    & Alanin Aminotransferaz (ALT) & 372.0 & 26.86 & 27.88 & 12.38 \\
318 & High Density Lipoprotein (HDL) & 273.0 & 49.48 & 17.72 & 12.62 \\
\hline
\end{tabular}
\caption{Values more than ten standard deviations from the feature mean.
Mean and standard deviation are computed on the uncorrected data.}
\label{tab:outliers}
\end{table}

Record 40 belongs to a patient without diabetes mellitus but reports a glucose
level of 575~mg/dl. Freckmann et al.~\cite{freckmann2007glucose} report a mean
24-hour interstitial glucose concentration of $89.3 \pm 6.2$~mg/dl in healthy
subjects under everyday life conditions. Even the highest postprandial values,
observed after breakfast, did not exceed 168~mg/dl. A value of 575~mg/dl is
therefore implausible for a non-diabetic patient, and the record was removed.

Record 239 reports an obesity percentage of 1954\%, which is physiologically
impossible. The height, weight and BMI of this record are mutually consistent
($65.9 / 1.50^2 \approx 29.3$). Obesity percentage is derived from BMI relative
to an ideal BMI, and a value of 19.54\% is consistent with the reference BMI
implied by the other records. We therefore attribute the value to a missing
decimal point and corrected it to 19.54\%.

Record 278 shows elevated levels of both AST (195~U/l) and ALT (372~U/l).
Although extreme, the simultaneous elevation of the two liver enzymes forms a
coherent clinical picture of liver involvement, which can occur in
gallstone disease when stones obstruct the bile ducts. These values were
therefore considered plausible and retained.

Record 318 reports an HDL level of 273~mg/dl. Since total cholesterol includes
both LDL and HDL, their sum cannot exceed it. In this record, however,
LDL + HDL $= 155 + 273 = 428$~mg/dl, exceeding the total cholesterol of
296~mg/dl. The value is therefore inconsistent, and the record was removed.

\paragraph{Features Cross-Correlation} 

Many bioimpedance features in the dataset measure closely related aspects of
body composition, so several of them carry nearly identical information.
Strongly correlated features add little predictive information, but they make
linear models harder to interpret: the coefficients become unstable, and the
effect of a single underlying quantity can be split between several features.
For this reason, we computed the pairwise Pearson correlation between all
features and removed redundant ones. Six pairs exceeded our threshold of
$|r| > 0.9$ (Table~\ref{tab:correlations}).

\begin{table}[h]
\centering
\begin{tabular}{llr}
\hline
Feature A & Feature B & $|r|$ \\
\hline
Total Body Fat Ratio (TBFR) (\%) & Lean Mass (LM) (\%) & 0.994 \\
Total Body Water (TBW) & Muscle Mass (MM) & 0.928 \\
Body Mass Index (BMI) & Obesity (\%) & 0.918 \\
Total Body Water (TBW) & Extracellular Water (ECW) & 0.904 \\
Intracellular Water (ICW) & Muscle Mass (MM) & 0.904 \\
Body Mass Index (BMI) & Total Fat Content (TFC) & 0.903 \\
\hline
\end{tabular}
\caption{Feature pairs with absolute Pearson correlation above 0.9.}
\label{tab:correlations}
\end{table}

\subsection{Models}

\subsubsection{Logistic Regression}

- What is Logistic Regression?
- Regularization parameter

\subsubsection{Logistic Regression + Log transformation}

- Why log transform some features

\subsubsection{Random Forest}

- What's Random Forest, How does it work?

\subsection{Performance Metric}

ROC AUC (Receiver Operating Characteristic - Area Under the Curve) is a standard evaluation metric for binary classification machine learning models.

\subsection{Feature evaluation}

Permutation feature importance is a model-agnostic technique that measures a feature's contribution by randomly shuffling its values and observing the resulting drop in model performance.

\section{Results}
\label{sec:results}



\section{Conclusion}
\label{sec:Conclusion}


\bibliographystyle{plainnat}
\bibliography{references}

\end{document}
